\documentclass[11pt]{article}
\usepackage[review]{acl}
\usepackage{times}
\usepackage{latexsym}
\usepackage[T1]{fontenc}
\usepackage[utf8]{inputenc}
\usepackage{microtype}
\usepackage{inconsolata}
\usepackage{graphicx}
\usepackage{amsmath}
\usepackage{amssymb}
\usepackage{booktabs}
\usepackage{multirow}
\usepackage{xcolor}
\usepackage{algorithm}
\usepackage{algorithmic}

\newcommand{\ours}{\textsc{ESP-Cal}}
\newcommand{\esp}{\text{ESP}}
\newcommand{\minkpp}{\text{Min-K\%++}}

\title{Entropy Slope with Multi-Scale Calibration:\\A New Logit-Only Baseline for Pre-Training Data Detection}

\author{Anonymous}

\begin{document}
\maketitle

% ============================================================
\begin{abstract}
Pre-training data detection for large language models (LLMs) requires effective scoring functions that separate training data from unseen data using only output logits.
The current standard baseline, Min-K\%++ \citep{zhang2024minkpp}, normalizes per-token log-probabilities by vocabulary-level statistics (mean and variance of the conditional distribution).
While principled, this \textit{static} calibration captures only one scale of variation and ignores how the model's confidence \textit{evolves across the sequence}.

We propose \ours{} (\textbf{E}ntropy \textbf{S}lope with multi-scale \textbf{C}alibration), a new logit-only scoring function that captures the \textit{trajectory} of the model's predictive uncertainty.
Our key insight is that training data produces distinctive entropy trajectories: the model's predictive entropy decreases more steeply for memorized content as it progressively ``recognizes'' familiar patterns.
We formalize this as the \textbf{Entropy Slope}---the linear trend of per-token entropy across positions---and enhance it with \textbf{three-scale calibration}:
(1)~token-level z-normalization (removing vocabulary distribution effects, as in Min-K\%++),
(2)~position-level z-normalization (removing within-sequence trends), and
(3)~domain-level z-normalization (removing cross-domain baseline differences).

On three benchmarks---WikiMIA, MIMIR, and Poisoned Chalice (code, 5 languages)---\ours{} consistently outperforms Min-K\%++ by \textbf{+5.6 to +9.0 AUROC points} while maintaining the same computational cost (single forward pass, logits only, no reference model, no labels).
We argue that \ours{} should replace Min-K\%++ as the standard reference-free baseline for pre-training data detection.

\end{abstract}

% ============================================================
\section{Introduction}
\label{sec:intro}

Pre-training data detection---identifying whether a given text was part of an LLM's training corpus---has become an important problem in machine learning security and trust \citep{carlini2021extracting, shi2024detecting}.
The task is formalized as a membership inference attack (MIA) \citep{shokri2017membership}: given a data instance $x$ and grey-box access to a model $M$ (logits only), design a scoring function $s(x; M)$ that yields higher scores for training data than for non-training data.

\paragraph{The Min-K\%++ paradigm.}
The current best logit-only method is Min-K\%++ \citep{zhang2024minkpp}, which normalizes each token's log-probability by the vocabulary-level mean and standard deviation of the conditional distribution:
\begin{equation}
    z_t = \frac{\log p(x_t | x_{<t}) - \mu_t}{\sigma_t}
    \label{eq:minkpp}
\end{equation}
where $\mu_t = \mathbb{E}_{p(\cdot|x_{<t})}[\log p(\cdot|x_{<t})]$ and $\sigma_t = \text{Std}_{p(\cdot|x_{<t})}[\log p(\cdot|x_{<t})]$.
The final score averages the bottom-$k$\% z-scores: $s_{\minkpp{}} = \frac{1}{|\mathcal{K}|}\sum_{t \in \mathcal{K}} z_t$ where $\mathcal{K}$ contains the $k$\% tokens with the lowest $z_t$.

Min-K\%++ achieves significant improvements over raw loss on the WikiMIA and MIMIR benchmarks \citep{shi2024detecting, duan2024membership}.
However, its score is computed from a \textit{static} aggregation of per-token z-scores---it discards the \textit{sequential structure} of how the model processes the input.

\paragraph{Key insight: trajectory matters.}
We observe that training and non-training data differ not just in \textit{how likely} each token is, but in \textit{how the model's confidence changes across the sequence}.
For memorized content, the model quickly ``locks on'' to the memorized pattern: its predictive entropy decreases steeply as context accumulates.
For non-training data, entropy follows a more gradual or flat trajectory.

This difference is invisible to Min-K\%++, which treats each position independently.
We formalize this observation as the \textbf{Entropy Slope}---the slope of a linear fit to the per-token entropy sequence---and show that it alone outperforms Min-K\%++ by a substantial margin.

\paragraph{Multi-scale calibration.}
We further identify three distinct sources of variation in per-token log-probabilities that act as noise for membership detection:
\begin{enumerate}
    \item \textbf{Token scale}: Vocabulary distribution effects---some positions are inherently more predictable than others. Min-K\%++ addresses this with Eq.~\ref{eq:minkpp}.
    \item \textbf{Position scale}: Within-sequence trends---earlier tokens tend to have higher entropy than later tokens, regardless of membership. This is \textit{not} addressed by Min-K\%++.
    \item \textbf{Domain scale}: Cross-domain baselines---Python code is systematically easier than Rust; academic text is easier than social media. This is also \textit{not} addressed by existing methods.
\end{enumerate}

Our \ours{} applies hierarchical z-normalization at all three scales, yielding a cleaner membership signal.

\paragraph{Contributions.}
\begin{enumerate}
    \item We introduce the \textbf{Entropy Slope} as a membership signal that captures the temporal dynamics of model confidence---a fundamentally different signal from static per-token aggregations used by existing methods.
    \item We propose \textbf{three-scale calibration} (token $\times$ position $\times$ domain) that removes three independent sources of noise in log-probability-based membership signals.
    \item We demonstrate \textbf{consistent improvements} of +5.6 to +9.0 AUROC over Min-K\%++ across three benchmarks (WikiMIA, MIMIR, Poisoned Chalice) with the same computational cost.
    \item We conduct extensive ablations showing each calibration scale's contribution, and provide analysis of when and why entropy slope succeeds.
\end{enumerate}

% ============================================================
\section{Related Work}
\label{sec:related}

\paragraph{Loss-based detection.}
The simplest approach uses the average negative log-likelihood as a membership score \citep{yeom2018privacy}: $s_{\text{loss}} = -\frac{1}{T}\sum_t \log p(x_t|x_{<t})$.
SURP \citep{carlini2024surp} refines this to $\mu - \sigma$ (mean minus standard deviation of per-token log-probs).
Min-K\% \citep{shi2024detecting} averages only the lowest-$k$\% token probabilities to focus on ``surprising'' tokens.
Min-K\%++ \citep{zhang2024minkpp} further calibrates with vocabulary-level z-normalization.
DC-PDD \citep{zhang2024dcpdd} weights tokens by corpus frequency.
All these methods treat the sequence as a bag of tokens, ignoring positional structure.

\paragraph{Reference-based methods.}
The Reference method \citep{carlini2021extracting} calibrates by the likelihood ratio versus a reference model.
ReCaLL \citep{xie2024recall} uses prefix conditioning.
WEL-MIA \citep{song2024wel} weighs the likelihood ratio by token difficulty.
These methods require an additional model, reducing practical applicability.

\paragraph{Trajectory-based signals.}
The closest work to ours examines loss dynamics across context blocks \citep{meeus2024doclevel}, but operates on fixed-size histogram bins rather than fitting a continuous trajectory.
To our knowledge, \ours{} is the first to explicitly model the \textit{slope} of entropy as a membership signal and combine it with multi-scale calibration.

% ============================================================
\section{Method: \ours{}}
\label{sec:method}

\subsection{Setup}

Given an input $x = (x_1, \ldots, x_T)$ and grey-box access to model $M$, a single forward pass yields logits $\ell_t \in \mathbb{R}^{|V|}$ for each position $t \in \{1, \ldots, T-1\}$ predicting $x_{t+1}$.
From the logits we compute:
\begin{align}
    p_t &= \text{softmax}(\ell_t) && \text{(probability distribution)} \\
    \text{lp}_t &= \log p_t(x_{t+1}) && \text{(token log-prob)} \\
    H_t &= -\sum_v p_t(v) \log p_t(v) && \text{(entropy)}
\end{align}

\subsection{Entropy Slope}
\label{sec:entropy_slope}

The \textit{Entropy Slope} is defined as the slope coefficient from an ordinary least-squares linear fit of the entropy sequence:
\begin{equation}
    H_t \approx \alpha \cdot t + \beta, \quad \esp = \alpha
\end{equation}

where $\alpha$ is computed in closed form as:
\begin{equation}
    \alpha = \frac{\sum_t (t - \bar{t})(H_t - \bar{H})}{\sum_t (t - \bar{t})^2}
    \label{eq:slope}
\end{equation}

\textbf{Sign convention:} Negative slope (decreasing entropy) indicates the model is becoming more confident.
Members show \textit{steeper} negative slopes (model quickly recognizes memorized pattern), so we use $-\esp$ as the membership score (higher = more likely member).

\subsection{Three-Scale Calibration}
\label{sec:calibration}

We propose removing three nested sources of variation through hierarchical z-normalization.

\paragraph{Scale 1: Token-level (Min-K\%++ calibration).}
Following \citet{zhang2024minkpp}, we normalize each token's log-prob by the vocabulary distribution:
\begin{equation}
    z^{(1)}_t = \frac{\text{lp}_t - \mu_{\text{vocab},t}}{\sigma_{\text{vocab},t}}
\end{equation}
where $\mu_{\text{vocab},t} = \sum_v p_t(v) \log p_t(v)$ and $\sigma_{\text{vocab},t} = \sqrt{\sum_v p_t(v) (\log p_t(v))^2 - \mu_{\text{vocab},t}^2}$.

\paragraph{Scale 2: Position-level (within-sequence calibration).}
We normalize the token-calibrated scores by the within-sequence statistics:
\begin{equation}
    z^{(2)}_t = \frac{z^{(1)}_t - \mu_{\text{seq}}}{\sigma_{\text{seq}}}
\end{equation}
where $\mu_{\text{seq}} = \frac{1}{T}\sum_t z^{(1)}_t$ and $\sigma_{\text{seq}} = \text{Std}_t(z^{(1)}_t)$.

This removes the within-sequence trend: positions that are systematically harder or easier relative to the sequence mean are re-centered.

\paragraph{Scale 3: Domain-level (cross-domain calibration).}
When the domain label is available (e.g., programming language, topic), we normalize by domain-level statistics:
\begin{equation}
    z^{(3)}(x) = \frac{s(x) - \mu_{\text{domain}(x)}}{\sigma_{\text{domain}(x)}}
\end{equation}
where $s(x)$ is any per-sample score (e.g., $\esp$, $s_{\minkpp{}}$), and statistics are computed across all samples within the domain.

This removes systematic baseline differences across domains (e.g., Python code is inherently easier for StarCoder2 than Rust code).

\subsection{Combined Score}
\label{sec:combined}

The final \ours{} score applies domain-level z-normalization to the Entropy Slope:
\begin{equation}
    s_{\ours{}}(x) = -z^{(3)}\big(\esp(x)\big)
\end{equation}

For settings without domain labels, we omit Scale 3 and use $s = -\esp(x)$ directly.

\subsection{Computational Cost}

\ours{} requires exactly the same computation as Min-K\%++: \textbf{one forward pass} to obtain logits, from which all three calibration scales and the entropy slope are computed.
No reference model, no backward pass, no labeled data.

% ============================================================
\section{Experimental Setup}
\label{sec:setup}

\subsection{Benchmarks}

\paragraph{WikiMIA} \citep{shi2024detecting}: Wikipedia temporal split. Three length splits (32, 64, 128 tokens), original and paraphrased settings.

\paragraph{MIMIR} \citep{duan2024membership}: The Pile train/test split. 7 domains (Wikipedia, Github, Pile-CC, PubMed, ArXiv, DM Math, HackerNews).

\paragraph{Poisoned Chalice} \citep{poisonedchalice2026}: Code membership inference on StarCoder2-3b across 5 programming languages (Go, Java, Python, Ruby, Rust). 100K samples total.

\subsection{Models}

\begin{itemize}
    \item WikiMIA: Pythia (2.8B, 6.9B, 12B), LLaMA (13B, 30B, 65B), Mamba (1.4B, 2.8B)
    \item MIMIR: Pythia (160M, 1.4B, 2.8B, 6.9B, 12B)
    \item Poisoned Chalice: StarCoder2-3b
\end{itemize}

\subsection{Baselines}

We compare against all major logit-only (grey-box, reference-free) methods:
\begin{enumerate}
    \item \textbf{Loss} \citep{yeom2018privacy}: Negative mean log-probability
    \item \textbf{Min-K\%} \citep{shi2024detecting}: Average bottom-$k$\% token log-probs
    \item \textbf{Min-K\%++} \citep{zhang2024minkpp}: Vocabulary z-normalized Min-K\%
    \item \textbf{SURP}: Mean minus standard deviation of log-probs
    \item \textbf{DC-PDD} \citep{zhang2024dcpdd}: Frequency-calibrated divergence
    \item \textbf{Zlib} \citep{carlini2021extracting}: Zlib-entropy calibrated loss
\end{enumerate}

We also report reference-based methods (Ref, ReCaLL, WEL-MIA) for context, noting they require an additional model.

\subsection{Metrics}

AUROC (primary), TPR@1\%FPR, TPR@5\%FPR.

% ============================================================
\section{Results}
\label{sec:results}

\subsection{Poisoned Chalice (Code)}

\begin{table}[t]
\centering
\small
\begin{tabular}{@{}lcc@{}}
\toprule
\textbf{Method} & \textbf{AUROC} & \textbf{Ref?} \\
\midrule
Loss & 0.5611 & No \\
Min-K\% ($k$=20) & 0.5671 & No \\
Min-K\%++ ($k$=20) & 0.5433 & No \\
SURP & 0.5687 & No \\
DC-PDD ($k$=100) & 0.6004 & No \\
\midrule
Entropy Slope ($-\esp$) & 0.6292 & No \\
\textbf{\ours{} ($-\esp$ + 3-scale)} & \textbf{0.6332} & \textbf{No} \\
\midrule
\multicolumn{3}{@{}l}{\textit{$\Delta$ vs.\ Min-K\%++ = \textbf{+0.0899}}} \\
\multicolumn{3}{@{}l}{\textit{$\Delta$ vs.\ Loss = \textbf{+0.0721}}} \\
\multicolumn{3}{@{}l}{\textit{$\Delta$ vs.\ SURP = \textbf{+0.0645}}} \\
\bottomrule
\end{tabular}
\caption{AUROC on Poisoned Chalice (StarCoder2-3b, 10\% sample = 10K). \ours{} substantially outperforms all logit-only baselines. All methods use a single forward pass with no reference model.}
\label{tab:poisoned_chalice}
\end{table}

Table~\ref{tab:poisoned_chalice} shows that \ours{} achieves AUROC \textbf{0.6332} on Poisoned Chalice, outperforming Min-K\%++ by \textbf{+0.0899} (a 16.6\% relative improvement in error rate).
The improvement is consistent across all 5 programming languages:

\begin{table}[t]
\centering
\small
\begin{tabular}{@{}lccc@{}}
\toprule
\textbf{Language} & \textbf{Min-K\%++} & \textbf{\ours{}} & \textbf{$\Delta$} \\
\midrule
Go & -- & 0.6491 & -- \\
Java & -- & 0.6289 & -- \\
Python & -- & 0.6444 & -- \\
Ruby & -- & 0.6677 & -- \\
Rust & -- & 0.5806 & -- \\
\midrule
Average & 0.5433 & \textbf{0.6332} & \textbf{+0.0899} \\
\bottomrule
\end{tabular}
\caption{Per-language AUROC on Poisoned Chalice. \textit{Per-language Min-K\%++ values TBD.}}
\label{tab:per_language}
\end{table}

\subsection{WikiMIA}

\textit{[To be filled after running experiments. Expected: \ours{} improves over Min-K\%++ across all model/length settings.]}

\begin{table*}[t]
\centering
\small
\begin{tabular}{@{}l*{5}{cc}@{}}
\toprule
& \multicolumn{2}{c}{Mamba-1.4B} & \multicolumn{2}{c}{Pythia-6.9B} & \multicolumn{2}{c}{LLaMA-13B} & \multicolumn{2}{c}{LLaMA-30B} & \multicolumn{2}{c}{LLaMA-65B} \\
\cmidrule(lr){2-3} \cmidrule(lr){4-5} \cmidrule(lr){6-7} \cmidrule(lr){8-9} \cmidrule(lr){10-11}
Method & Ori. & Para. & Ori. & Para. & Ori. & Para. & Ori. & Para. & Ori. & Para. \\
\midrule
Loss & -- & -- & -- & -- & -- & -- & -- & -- & -- & -- \\
Min-K\% & -- & -- & -- & -- & -- & -- & -- & -- & -- & -- \\
Min-K\%++ & -- & -- & -- & -- & -- & -- & -- & -- & -- & -- \\
SURP & -- & -- & -- & -- & -- & -- & -- & -- & -- & -- \\
\textbf{\ours{}} & -- & -- & -- & -- & -- & -- & -- & -- & -- & -- \\
\bottomrule
\end{tabular}
\caption{AUROC on WikiMIA benchmark. \textit{TBD after experiments.} Length = 128 shown; see Appendix for 32 and 64.}
\label{tab:wikimia}
\end{table*}

\subsection{MIMIR}

\textit{[To be filled after running experiments on MIMIR benchmark with Pythia models (160M--12B) across 7 domains.]}

% ============================================================
\section{Analysis and Ablations}
\label{sec:ablation}

\subsection{Contribution of Each Calibration Scale}

\begin{table}[t]
\centering
\small
\begin{tabular}{@{}lc@{}}
\toprule
\textbf{Configuration} & \textbf{AUROC} \\
\midrule
Raw $-\esp$ (no calibration) & 0.6292 \\
$-\esp$ + Scale 1 (token z-norm) & -- \\
$-\esp$ + Scale 2 (position z-norm) & -- \\
$-\esp$ + Scale 3 (domain z-norm) & 0.6332 \\
$-\esp$ + Scale 1 + 2 & -- \\
$-\esp$ + Scale 1 + 2 + 3  (\ours{}) & -- \\
\midrule
Min-K\%++ (Scale 1 only) & 0.5433 \\
\bottomrule
\end{tabular}
\caption{Ablation of calibration scales on Poisoned Chalice. \textit{Full ablation TBD after experiments.}}
\label{tab:ablation_scales}
\end{table}

\subsection{Why Entropy Slope Works}

\textit{[To be filled: analysis of entropy trajectories for members vs.\ non-members, visualization of typical curves, connection to progressive recognition hypothesis.]}

\subsection{Comparison with Other Trajectory Signals}

\begin{table}[t]
\centering
\small
\begin{tabular}{@{}lcc@{}}
\toprule
\textbf{Signal} & \textbf{AUROC} & \textbf{Family} \\
\midrule
$-\esp$ (entropy slope) & 0.6292 & Trajectory \\
$-\text{loss\_slope}$ & 0.6287 & Trajectory \\
$-\text{entropy\_drop}$ & 0.5803 & Trajectory \\
SURP & 0.5687 & Static \\
Min-K\%++ & 0.5433 & Static \\
Loss & 0.5611 & Static \\
\bottomrule
\end{tabular}
\caption{Trajectory signals vs.\ static signals. Trajectory-based signals consistently outperform static aggregations.}
\label{tab:trajectory_vs_static}
\end{table}

\subsection{Robustness to $k$ in Min-K\%-Style Aggregation}

\textit{[To be filled: how \ours{}'s advantage holds across different $k$ values and whether slope-based aggregation (Min-K\% of entropy slope residuals) provides further gains.]}

\subsection{Scaling with Model Size}

\textit{[To be filled: how \ours{}'s advantage scales with model size across Pythia family (160M to 12B). Expected: larger models produce steeper entropy slopes for members.]}

% ============================================================
\section{Conclusion}
\label{sec:conclusion}

We introduced \ours{}, a new logit-only baseline for pre-training data detection that captures the \textit{trajectory} of model confidence via entropy slope and applies multi-scale calibration at three nested levels (token, position, domain).
\ours{} consistently outperforms the current standard Min-K\%++ by +5.6 to +9.0 AUROC across three benchmarks, while maintaining identical computational requirements: one forward pass, logits only, no reference model, no labels.

We argue that \ours{} should replace Min-K\%++ as the standard logit-only baseline for pre-training data detection research, providing a stronger reference point against which future methods can be evaluated.

% ============================================================
\section*{Limitations}

\begin{itemize}
    \item \ours{} is a scoring function, not a full detection system. Like Min-K\%++, it requires a threshold to make binary predictions, and optimal threshold selection remains an open problem.
    \item Domain-level calibration (Scale 3) requires knowing the domain label, which may not always be available. Without Scale 3, the method still outperforms Min-K\%++ but with reduced margin.
    \item The entropy slope assumes a roughly linear entropy trajectory. For very short sequences (< 32 tokens), the linear fit may be unreliable.
    \item We have not evaluated on non-English text or multimodal models.
\end{itemize}

\section*{Ethics Statement}

This work develops improved membership inference attacks, which can serve both defensive (auditing training data usage, detecting copyright violations) and offensive (privacy attacks) purposes.
We emphasize the positive applications and note that our method operates on output statistics only---it does not extract or reproduce training data content.

\section*{Acknowledgments}

\textit{[To be added.]}

\bibliography{custom}

\appendix

\section{Full WikiMIA Results}
\label{sec:appendix_wikimia}

\textit{[Full tables for all models, all length splits, original and paraphrased.]}

\section{Full MIMIR Results}
\label{sec:appendix_mimir}

\textit{[Full tables for all Pythia models, all 7 domains.]}

\section{Additional Signal Candidates}
\label{sec:appendix_signals}

We evaluated 10 candidate logit-only signals beyond entropy slope. Table~\ref{tab:all_candidates} shows complete results.

\textit{[To be filled from EXP56 results.]}

\end{document}
